\documentclass[11pt]{article}

\usepackage[a4paper,margin=1in]{geometry}
\usepackage{amsmath,amssymb,bm}
\usepackage{physics}
\usepackage{mathtools}
\usepackage{enumitem}
\usepackage{microtype}
\usepackage[hidelinks]{hyperref}

\title{Origin of the Positive- and Negative-Frequency Convention}
\author{}
\date{}

\begin{document}

\maketitle

\section{The central question}

Why is a factor of the form
\begin{equation}
e^{-i\omega t},
\qquad \omega>0,
\label{eq:exp-minus}
\end{equation}
called a \emph{positive-frequency} component, even though the exponent contains a minus sign?

The answer is that the terms ``positive frequency'' and ``negative frequency'' refer to the sign of the Fourier-frequency variable, not to the explicit sign multiplying $i\omega t$ in isolation. Which exponential represents positive frequency therefore depends on the Fourier-transform convention being used.

\section{Fourier-transform convention}

A common convention in physics is
\begin{equation}
f(t)
=
\int_{-\infty}^{\infty}
\frac{d\omega}{2\pi}\,
\widetilde f(\omega)e^{-i\omega t},
\label{eq:inverse-fourier}
\end{equation}
with the corresponding forward transform
\begin{equation}
\widetilde f(\omega)
=
\int_{-\infty}^{\infty}
dt\,
f(t)e^{+i\omega t}.
\label{eq:forward-fourier}
\end{equation}

In Eq.~\eqref{eq:inverse-fourier}, the functions
\begin{equation}
e^{-i\omega t}
\label{eq:fourier-basis}
\end{equation}
are the Fourier basis functions, and $\omega$ is the variable that labels their frequency.

Therefore, if
\begin{equation}
\omega>0,
\end{equation}
then
\begin{equation}
e^{-i\omega t}
\end{equation}
is, by definition within this Fourier convention, a \emph{positive-frequency Fourier component}.

Conversely,
\begin{equation}
e^{+i\omega t}
=
e^{-i(-\omega)t}.
\label{eq:rewrite-positive-exp}
\end{equation}
For $\omega>0$, this is a Fourier component whose frequency label is $-\omega<0$. It is therefore called a \emph{negative-frequency component}.

Thus, under the convention of Eq.~\eqref{eq:inverse-fourier},
\begin{equation}
\boxed{
\begin{aligned}
e^{-i\omega t},\quad \omega>0
&\quad\longleftrightarrow\quad
\text{positive frequency},\\[4pt]
e^{+i\omega t},\quad \omega>0
&\quad\longleftrightarrow\quad
\text{negative frequency}.
\end{aligned}}
\label{eq:frequency-identification}
\end{equation}

\section{Explicit Fourier-transform examples}

Consider first
\begin{equation}
f(t)=e^{-i\omega_0 t},
\qquad
\omega_0>0.
\label{eq:positive-example}
\end{equation}
Using Eq.~\eqref{eq:forward-fourier},
\begin{align}
\widetilde f(\omega)
&=
\int_{-\infty}^{\infty}
dt\,
e^{-i\omega_0t}e^{+i\omega t}\\
&=
\int_{-\infty}^{\infty}
dt\,
e^{i(\omega-\omega_0)t}\\
&=
2\pi\delta(\omega-\omega_0).
\label{eq:positive-spectrum}
\end{align}

The spectrum is concentrated at
\begin{equation}
\omega=+\omega_0.
\end{equation}
Hence $e^{-i\omega_0t}$ is a positive-frequency component.

Now consider
\begin{equation}
g(t)=e^{+i\omega_0t}.
\label{eq:negative-example}
\end{equation}
Its Fourier transform is
\begin{align}
\widetilde g(\omega)
&=
\int_{-\infty}^{\infty}
dt\,
e^{+i\omega_0t}e^{+i\omega t}\\
&=
\int_{-\infty}^{\infty}
dt\,
e^{i(\omega+\omega_0)t}\\
&=
2\pi\delta(\omega+\omega_0).
\label{eq:negative-spectrum}
\end{align}

The spectrum is therefore concentrated at
\begin{equation}
\omega=-\omega_0.
\end{equation}
Thus $e^{+i\omega_0t}$ is a negative-frequency component.

\section{Why the terminology can initially look backwards}

The expression
\begin{equation}
e^{-i\omega t}
\end{equation}
contains a minus sign, so it is tempting to associate it with ``negative frequency.'' However, that minus sign belongs to the chosen Fourier kernel.

The standard basis has been chosen as
\begin{equation}
e^{-i\omega t}.
\end{equation}
The sign of the frequency is determined by the value of the label $\omega$.

A useful way to avoid confusion is always to rewrite a time-dependent exponential in the standard Fourier form
\begin{equation}
e^{-i\Omega t}.
\label{eq:standard-form}
\end{equation}
Then inspect the sign of $\Omega$.

For example,
\begin{equation}
e^{-i5t}
\end{equation}
has $\Omega=+5$ and is therefore positive frequency, whereas
\begin{equation}
e^{+i5t}
=
e^{-i(-5)t}
\end{equation}
has $\Omega=-5$ and is therefore negative frequency.

\section{A real oscillation contains both signs of frequency}

Consider a real monochromatic field,
\begin{equation}
E(t)=E_0\cos(\omega_0t),
\qquad
\omega_0>0.
\label{eq:real-field}
\end{equation}
Using Euler's identity,
\begin{equation}
\cos(\omega_0t)
=
\frac{1}{2}
\left(
e^{-i\omega_0t}
+
e^{+i\omega_0t}
\right),
\label{eq:cos-decomposition}
\end{equation}
so
\begin{equation}
E(t)
=
\frac{E_0}{2}e^{-i\omega_0t}
+
\frac{E_0}{2}e^{+i\omega_0t}.
\end{equation}

We therefore define
\begin{equation}
E(t)=E^{(+)}(t)+E^{(-)}(t),
\label{eq:E-decomposition}
\end{equation}
where
\begin{align}
E^{(+)}(t)
&=
\frac{E_0}{2}e^{-i\omega_0t},
\label{eq:E-positive}\\
E^{(-)}(t)
&=
\frac{E_0}{2}e^{+i\omega_0t}.
\label{eq:E-negative}
\end{align}

For a real field,
\begin{equation}
E^{(-)}(t)
=
\left[E^{(+)}(t)\right]^*.
\label{eq:conjugate-relation}
\end{equation}

Thus a real sinusoidal field necessarily contains equal positive- and negative-frequency contributions.

\section{The convention could have been chosen oppositely}

The terminology is not an intrinsic property of the exponential alone. Suppose instead that one adopted the inverse Fourier transform
\begin{equation}
f(t)
=
\int_{-\infty}^{\infty}
\frac{d\omega}{2\pi}\,
\widetilde f(\omega)e^{+i\omega t}.
\label{eq:opposite-convention}
\end{equation}

With this convention, the basis function associated with a positive value $\omega>0$ would be
\begin{equation}
e^{+i\omega t}.
\end{equation}

Therefore the assignment
\begin{equation}
e^{-i\omega t}
\longleftrightarrow
\text{positive frequency}
\end{equation}
is not a convention-independent mathematical truth. It follows from the Fourier-sign convention in Eq.~\eqref{eq:inverse-fourier}.

What is physically important is consistency: once a Fourier convention is chosen, all definitions of positive- and negative-frequency components must follow that convention.

\section{Why the $e^{-i\omega t}$ convention is natural in quantum mechanics}

The convention commonly used in atomic physics and quantum optics is particularly convenient because it agrees naturally with Schr\"odinger time evolution.

Let $\ket{n}$ be an energy eigenstate satisfying
\begin{equation}
H\ket{n}=E_n\ket{n}.
\label{eq:energy-eigenstate}
\end{equation}
Its Schr\"odinger-picture time dependence is
\begin{equation}
\ket{n(t)}
=
e^{-iE_nt/\hbar}\ket{n}.
\label{eq:schrodinger-evolution}
\end{equation}

Defining
\begin{equation}
\omega_n=\frac{E_n}{\hbar},
\label{eq:omega-energy}
\end{equation}
we obtain
\begin{equation}
\ket{n(t)}
=
e^{-i\omega_nt}\ket{n}.
\label{eq:energy-frequency-evolution}
\end{equation}

For positive energy relative to the chosen reference,
\begin{equation}
E_n>0
\quad\Longrightarrow\quad
\omega_n>0,
\end{equation}
and the corresponding time evolution contains $e^{-i\omega_nt}$.

This makes the Fourier convention
\begin{equation}
f(t)\sim e^{-i\omega t}
\end{equation}
especially convenient in quantum mechanics.

\section{Application to the electromagnetic field}

In quantum optics, the electric field is decomposed as
\begin{equation}
\widehat{\bm E}(\bm r,t)
=
\widehat{\bm E}^{(+)}(\bm r,t)
+
\widehat{\bm E}^{(-)}(\bm r,t).
\label{eq:quantum-field-decomposition}
\end{equation}

For a single plane-wave mode, schematically,
\begin{align}
\widehat{\bm E}^{(+)}(\bm r,t)
&\propto
\bm\epsilon\,
\hat a\,
e^{i\bm k\cdot\bm r}
e^{-i\omega t},
\label{eq:quantum-positive-field}\\
\widehat{\bm E}^{(-)}(\bm r,t)
&\propto
\bm\epsilon^*\,
\hat a^\dagger\,
e^{-i\bm k\cdot\bm r}
e^{+i\omega t}.
\label{eq:quantum-negative-field}
\end{align}

Here:
\begin{itemize}[leftmargin=2em]
    \item $\hat a$ is the photon-annihilation operator,
    \item $\hat a^\dagger$ is the photon-creation operator,
    \item $\bm\epsilon$ is the polarization vector,
    \item $\bm k$ is the wave vector.
\end{itemize}

The superscripts $(+)$ and $(-)$ in
$\widehat{\bm E}^{(+)}$ and $\widehat{\bm E}^{(-)}$
refer to the sign of the Fourier frequency, not to the sign appearing explicitly in the exponential.

\section{Connection to resonant light--atom interaction}

The convention is particularly useful when comparing the optical field with the time dependence of atomic transition operators.

Let
\begin{equation}
\omega_{eg}
=
\frac{E_e-E_g}{\hbar}
>0
\label{eq:transition-frequency}
\end{equation}
be an atomic transition frequency.

The atomic raising operator
\begin{equation}
\ket{e}\bra{g}
\end{equation}
acquires the interaction-picture time dependence
\begin{equation}
\ket{e}\bra{g}
\longrightarrow
e^{+i\omega_{eg}t}
\ket{e}\bra{g}.
\label{eq:raising-time}
\end{equation}

The positive-frequency optical field contains
\begin{equation}
E^{(+)}(t)\propto e^{-i\omega t}.
\label{eq:optical-positive-time}
\end{equation}

Their product therefore varies as
\begin{equation}
e^{+i\omega_{eg}t}
e^{-i\omega t}
=
e^{i(\omega_{eg}-\omega)t}.
\label{eq:near-resonant-factor}
\end{equation}

When $\omega\simeq\omega_{eg}$, this factor varies slowly. This is the near-resonant term retained in the rotating-wave approximation.

By contrast, combining the atomic raising operator with the negative-frequency field gives
\begin{equation}
e^{+i\omega_{eg}t}
e^{+i\omega t}
=
e^{i(\omega_{eg}+\omega)t},
\label{eq:counterrotating-factor}
\end{equation}
which oscillates rapidly near resonance and is the corresponding counter-rotating term.

Thus the positive-frequency convention is not merely convenient bookkeeping: once adopted consistently, it makes the resonant structure of the light--atom interaction transparent.

\section{Summary}

Under the commonly used Fourier convention
\begin{equation}
f(t)
=
\int\frac{d\omega}{2\pi}
\widetilde f(\omega)e^{-i\omega t},
\end{equation}
the exponential $e^{-i\omega t}$ with $\omega>0$ is called a positive-frequency component because its Fourier spectrum is located at positive $\omega$.

Similarly,
\begin{equation}
e^{+i\omega t}
=
e^{-i(-\omega)t}
\end{equation}
has Fourier frequency $-\omega$ and is therefore a negative-frequency component.

The essential statement is
\begin{equation}
\boxed{
\text{The signs ``positive'' and ``negative'' refer to the Fourier-frequency label,
not directly to the sign in the exponent.}
}
\label{eq:main-summary}
\end{equation}

The assignment depends on the Fourier-transform convention, but the convention using $e^{-i\omega t}$ for positive frequency is especially natural in quantum mechanics because positive-energy stationary states evolve as $e^{-iEt/\hbar}$.

\end{document}
